\paragraph{AQ-136-FIRST-SKELETON: a common free integral first skeleton by actual replacements.}
The supplied intake reports a common first-skeleton replacement
as A-004. Its prior full proof is unavailable here. We give a
complete construction using only the original uniform
admissibility tests. This is a reduction of the integral
comparison problem, not a proof of integral connectedness.

Let $k$ be any commutative ring and let
$C_1,\ldots,C_t\in\mathcal A_k$ be a finite nonempty family.
There is an integer $N\ge1$, one $1$-truncated simplicial strict
PD differential $S_k$-algebra $F$, and admissible theories
$C'_j$ with strict maps $C'_j\to C_j$, such that every $C'_j$
has exactly the same first skeleton $F$. Explicitly,
\[
 \begin{gathered}
 F_0=S_k,\qquad
 F_1=S_k\otimes_k
       \Gamma_k(x_m:m\ge N)\otimes_k\Lambda_k(e_m:m\ge N),\\
 |x_m|=(m,0),\quad |e_m|=(m,1),\quad dx_m=e_m,\quad de_m=0,\\
 d_0x_m=0,\quad d_1x_m=T^{[m]},\quad
 d_0e_m=d_1e_m=0,
 \end{gathered}                                             \tag{136.SK1}
\]
with the degeneracy $F_0\to F_1$ given by coefficients.
Interchanging the two endpoint names makes no difference.
The augmentation $S_k\to C'_{j,n}$ still has the original
uniform vertical cohomology comparison, and every horizontal
row is acyclic in its weight tail.

\begin{proof}
\emph{1. Choose the common first skeleton.}
Choose $N$ so that every zero-form row $(C_j)_*^{m,0}$ is
acyclic whenever $m\ge N$. In normalized degree zero its
coefficient $T^{[m]}$ is therefore a boundary. Choose
\[
 y_{j,m}\in (C_j)_1^{m,0},\qquad
       d_0y_{j,m}=0,\quad d_1y_{j,m}=T^{[m]}.
                                                               \tag{136.SK2}
\]
One chooses the sign of a normalized primitive to obtain
these two endpoints, independently of the normalization's
boundary-sign convention. The assignments
$x_m\mapsto y_{j,m}$ and $e_m\mapsto dy_{j,m}$, together with
coefficients in degree zero, define a strict map from $F$
to the first truncation of $C_j$. The free PD and exterior
universal properties verify the algebra and differential
relations, and the two faces and sole degeneracy verify
all truncated simplicial identities.

Every vertical complex $F_1^{w,*}$ has coefficient cohomology
$S_k^w$ in degree zero and zero in higher degrees, in every
weight. This is the integral PD Poincare contraction of
independent even zero-form disks. It is valid for the
infinitely many displayed variables by taking the filtered
union over finite sets; at a fixed weight only finitely many
of their weights can occur. The contraction is $S_k$-linear
and weight-preserving.

\emph{2. Two facts about the first coskeleton.}
For any $1$-truncated simplicial algebra $E$, let
$\operatorname{cosk}_1E$ denote the right Kan extension to
all simplicial degrees, computed by the limits of its
vertex and edge data. Its value at $n$ consists of compatible
sections on the graph $G_n=\operatorname{sk}_1\Delta[n]$:
one value in $E_0$ at each vertex and one in $E_1$ at each
edge, with the specified endpoints. Degenerate edges add
no new variables. Operations are computed componentwise.
This is the finite-limit formula of the Stacks Project,
Lemmas~14.19.2--14.19.4 in Section~14.19, Tag~0AMA;
in particular the first truncation is unchanged.

First, if $M$ is a $1$-truncated simplicial module, then the
normalization of its coskeleton is
\[
 \cdots\longrightarrow0\longrightarrow
       \ker(\partial:N_1M\to M_0)
       \xrightarrow{\;\mathrm{incl}\;}N_1M
       \xrightarrow{\partial}M_0.
                                                               \tag{136.SK3}
\]
The displayed kernel is in degree two. One direct verification
uses triangle edge values: imposing the first two normalized
face conditions leaves exactly one edge value with both
endpoints zero, and the last face gives its inclusion in
$N_1M$. In degree at least three, every edge and vertex lies
in one of the normalized zero faces, so the normalized term
is zero. Thus the only possible horizontal homology of this
coskeleton is $H_0=\operatorname{coker}\partial$.

Second, suppose in a fixed weight the map
$E_1^q\to E_0^q\oplus E_0^q$ is surjective in every degree
needed for a cohomology calculation. Then sections on
$G_n$ are the kernel of the exact incidence sequence
\[
 0\longrightarrow(\operatorname{cosk}_1E)_n
 \longrightarrow E_0^{\,V(G_n)}\oplus E_1^{\,E(G_n)}
 \longrightarrow (E_0\oplus E_0)^{E(G_n)}
 \longrightarrow0.                                      \tag{136.SK4}
\]
The last map subtracts each edge's endpoint values from the
values at its two incident vertices. It is surjective by
choosing the two endpoint values independently on each edge.
If the cohomology of both $E_0,E_1$ is canonically $S_k^w$
in degree zero and zero above, the resulting cohomology of
the kernel is exactly
\[
                 H^r(G_n;k)\otimes_k S_k^w.                \tag{136.SK5}
\]
Indeed the degree-zero incidence complex in (136.SK4) is
the vertex-edge cochain complex of the graph after eliminating
one copy of the scalar for each edge. It has $H^0=S_k^w$
because $G_n$ is connected, its usual graph $H^1$, and no
higher cohomology. All identifications are natural in $E$.
If the hypotheses hold only through degree $d+1$, the same
image-complex truncation as in AQ-136-COSPAN proves (136.SK5)
through degree $d$. Finite products over a graph are finite
direct sums of modules, so no product exactness hypothesis is
needed. The cutoff does not depend on the number of vertices
or edges.

\emph{3. The two coskeleta have the same vertical tail comparison.}
Put $U=\operatorname{cosk}_1F$ and
$V_j=\operatorname{cosk}_1(\operatorname{sk}_1 C_j)$.
The maps of Step~1 induce $U\to V_j$.
For $F$, endpoint surjectivity holds in zero forms for all
$w\ge N$: the coefficient $T^{[w]}$ and the generator
$x_w$ give the two independent endpoint values. In positive
form degrees the endpoint target is zero, since $F_0=S_k$.
Step~1 and (136.SK4)--(136.SK5) therefore compute the vertical
cohomology of $U_n$ exactly as graph cohomology for $w\ge N$.

For $C_j$, acyclicity of its form-$q$ row gives the uniform
endpoint surjectivity $(C_j)_1^{w,q}\to
(C_j)_0^{w,q}\oplus(C_j)_0^{w,q}$ by (136.CS2), with $n=0$.
Its uniform vertical comparison in simplicial degrees zero
and one and the finite-degree version of (136.SK4) give the
same graph-cohomology answer for $V_{j,n}$, in any specified
finite range of form degrees, once $w$ is sufficiently large.
These bounds are independent of $n$.
The map $U\to V_j$ preserves coefficients and hence is the
identity on the incidence cohomology calculation. Consequently
\[
      H^q(U_n^{w,*})\xrightarrow{\sim} H^q(V_{j,n}^{w,*})
                                                               \tag{136.SK6}
\]
holds in every fixed finite range, uniformly in $n$, in a
weight tail. Neither coskeleton is claimed to be admissible:
its graph $H^1$ usually remains nonzero.

\emph{4. The actual replacements and their horizontal test.}
Define the strict PD pullback
\[
                 C'_j=C_j\times_{V_j}U,                    \tag{136.SK7}
\]
using the unit $C_j\to V_j$ of truncation and coskeleton.
For each fixed form degree $q$ and sufficiently large weight,
this unit is degreewise surjective as a map of simplicial
modules. To see this, a compatible family on $G_n$ is a
chain map from $N(k[G_n])$ into the acyclic normalized row
of $C_j$. Extend it across $\Delta[n]$ by the same increasing
degree filling argument as in (136.CS2). The relative
normalized chains have a free basis, and all relative cells
have degree at least two. This proves surjectivity uniformly
in $n$ from the one row cutoff.
Thus in those weights (136.SK7) gives an exact module sequence
\[
       0\longrightarrow (C'_j)_*^{w,q}
       \longrightarrow (C_j)_*^{w,q}\oplus U_*^{w,q}
       \longrightarrow (V_j)_*^{w,q}\longrightarrow0.
                                                               \tag{136.SK8}
\]

The rows of $U$ are acyclic in a uniform weight tail by
(136.SK3). In form degree zero the image of $\partial$
contains $kT^{[w]}$ for every $w\ge N$, by the generator
$x_w$; in positive form degree its degree-zero target is
already zero. Likewise the rows of $V_j$ are acyclic in
the same tails as those of $C_j$, since (136.SK3) has only
the original row's $H_0$. The normalized exact sequence of
(136.SK8) proves horizontal acyclicity for $C'_j$.

\emph{5. The vertical test and the exact first skeleton.}
Fix $d\ge1$. Choose $w$ large enough for (136.SK8) to be
degreewise exact through degree $d+1$, for the coefficient
comparison of $C_j$ through that range, and for (136.SK6)
through that range. The finite maximum of the required
cutoffs is independent of simplex dimension.
Let $I$ be the image of $C_j\oplus U\to V_j$ as a vertical
complex in a fixed simplex and weight. It agrees with $V_j$
through degree $d+1$, so cohomology through degree $d$ is
unchanged. In the long exact sequence of
$0\to C'_j\to C_j\oplus U\to I\to0$, the map contributed
by $U$ is the isomorphism (136.SK6). In degree zero the
matrix is the scalar difference $S_k^w\oplus S_k^w\to S_k^w$,
with diagonal kernel. It follows that $C'_j$ has coefficient
$H^0=S_k^w$ and zero $H^q$ for $1\le q\le d$.
The maps and coefficients in (136.SK7) are strict PD maps,
so this proves $C'_j\in\mathcal A_k$.

Finally the unit $C_j\to\operatorname{cosk}_1\operatorname{sk}_1C_j$
is an isomorphism in simplicial degrees zero and one.
The pullback in those degrees is therefore exactly $F$,
with its displayed faces, degeneracy, products, divided
powers and differential. Thus all the replacements have
the same first skeleton, and the projection in (136.SK7)
is the asserted admissible arrow to the original target.
\end{proof}

The theorem gives actual arrows, so replacing targets in
this way does not change their connected components in the
ordinary all-arrow localization. It does not assert that
arbitrary maps between the replacements fix their common
first skeleton, or that a later prescribed partial model
always has a completion. Higher-simplex compatibility,
including integral secondary periods, remains to be established.
